\section{The Craft of Dying}

In the fifteenth century the Church's long experience at the deathbed was gathered into a short Latin treatise on the art of dying. Its author is unknown; it cites ``the Chancellor of Paris,'' Jean Gerson, among its authorities. Many printed versions of it exist, and the picture-books of the same art, its editor records, are found ``in almost every European language.''\endnote{Frances M. M. Comper, ``Note on the Book of the Craft of Dying,'' in \emph{The Book of the Craft of Dying}, ed. Comper (1917), 48--50.} In English it was called \emph{The Book of the Craft of Dying}, and it was written for everyone: ``Forasmuch as the passage of death, of the wretchedness of the exile of this world, for uncunning of dying---not only to lewd men but also to religious and devout persons---seemeth wonderfully hard and perilous, and also right fearful and horrible,'' the book offers ``a short manner of exhortation, for teaching and comforting of them that be in point of death.''\endnote{\emph{The Book of the Craft of Dying}, prologue, ed. Comper (1917), 3.} The ``lewd men'' are laymen, as the editor explains, and ``uncunning'' is ignorance. It is profitable, the book says, ``to all true Christian men, to learn and have craft and knowledge to die well.''

The book begins with the commendation of death, and opens it with a distinction. Bodily death is dreadful, ``yet spiritual death of the soul is as much more horrible and detestable, as the soul is more worthy and precious than the body.'' The death of the good man, on the other hand, ``is ever precious in the sight of God, what manner of bodily death that ever they die.''\endnote{\emph{Craft of Dying}, ch.~1, Comper, 5; with Psalm 115:15, ``Precious in the sight of the Lord is the death of his saints,'' Douay--Rheims (Challoner).} For such a man, it says, death is ``a going [out] of prison, and an ending of exile; a discharging of an heavy burden, that is the body; finishing of all infirmities; a scaping of all perils; destroying of all evil things; breaking of all bonds; paying of [the] debt of natural duty; turning again into his country; and entering into bliss and joy.''\endnote{\emph{Craft of Dying}, ch.~1, Comper, 6, brackets as printed.} Scripture uses the same images. The Psalmist prays, ``Bring my soul out of prison, that I may praise thy name''; the Wise Man says that ``the corruptible body is a load upon the soul''; and the Apostle that ``if our earthly house of this habitation be dissolved,'' we have ``a house not made with hands, eternal in heaven.''\endnote{Psalm 141:8; Wisdom 9:15; 2 Corinthians 5:1; Douay--Rheims (Challoner).} Then the book defines dying well in a sentence it borrows from a wise man: ``To die well is to die gladly and wilfully,'' that is, to give up to God with a willing heart the life that must be given up in any case.\endnote{\emph{Craft of Dying}, ch.~1, Comper, 7.}

\subsection{The five temptations}

The second chapter warns that the dying ``have greatest and most grievous temptations, and such as they never had before in all their life,'' and it names five.\endnote{\emph{Craft of Dying}, ch.~2, Comper, 9--10.} The devil attacks the faith, ``forasmuch as faith is fundament of all men's soul's-heal.'' He urges despair, setting a man's sins before him. He stirs impatience in the sufferer. He tempts the devout most subtly of all, with complacence, ``that is spiritual pride, with the which the devil tempteth and beguileth most religious, and devout and perfect men,'' whispering, ``O how stable art thou in the faith! how strong in hope! how sad in patience! O how many good deeds hast thou done!'' And he holds worldly men by ``their wives, their children, their carnal friends, and their worldly riches, and other things that they have loved inordinately before.''\endnote{\emph{Craft of Dying}, ch.~2, Comper, 10, 18--20.} Against all five the book sets one consolation. The devil ``may compel no man, nor in no manner of wise prevail against him for to consent to him---as long as a man hath the use of reason with him---but if he will wilfully consent unto him.''\endnote{\emph{Craft of Dying}, ch.~2, Comper, 20--21.} God, as the book renders the Apostle, ``is true, and will not suffer you to be tempted more than ye may bear.''

\subsection{The questions at the bedside}

Then the book gives those who sit with the dying a set of questions to ask, which it attributes to Saint Anselm. Each is short and asks only for a yes:

\begin{quote}
Brother, art thou glad that thou shalt die in the faith of Christ? The sick man answereth: Yea.

Knowest thou well that thou hast not done as thou shouldst have done? He answereth: Yea.

Repentest thee thereof? He answereth: Yea.

Hast thou full will to amend thee, if thou mightest have full space of life? He answereth: Yea.

Believest thou fully that Our Lord Jesu Christ, God's Son, died for thee? He sayeth: Yea.

Thankest thou Him thereof with all thine heart? He answereth: Yea.

Believest thou verily that thou mayest not be saved but by Christ's [death and His] passion? He answereth: Yea.\endnote{\emph{Craft of Dying}, ch.~3, Comper, 22--23, brackets as printed.}
\end{quote}

\noindent After the questions the book puts into the dying man's mouth, against the terror of the judgment, a prayer of great simplicity: ``Lord, I put the death of Our Lord Jesu Christ between me and mine evil deeds, between me and the judgment; otherwise will I not strive with Thee.'' If God should say to him that he has deserved damnation, he is to answer, ``The death of our Lord Jesu Christ I put between me and mine evil merits.'' Last of all he is to say three times, \latin{In manus tuas, Domine}, and if he can no longer speak, those who stand about him are to say it for him: \latin{In manus tuas, Domine, commendamus spiritum eius}.\endnote{\emph{Craft of Dying}, ch.~3, Comper, 23.}

The fifth chapter turns to the living and rebukes them. ``It is greatly to be noted, and to be taken heed of, that right seldom (that) any man---yea among religious and devout men---dispose themselves to death betimes as they ought. For every man weeneth himself to live long, and troweth not that he shall die in short time; and doubtless that cometh of the devil's subtle temptation.''\endnote{\emph{Craft of Dying}, ch.~5, Comper, 32.} Men seek medicine for the body sooner than for the soul. The sick are flattered with hopes of recovery, and ``as the great Clerk, the Chancellor of Paris saith: Often times by such a [vain and a] false cheering and comforting, and feigned behoting of bodily heal, and trusting thereupon, men run and fall into certain damnation everlastingly.''\endnote{\emph{Craft of Dying}, ch.~5, Comper, 33.} The ``behoting'' is the promising. Those at the bedside should therefore tell the sick the truth while there is still time for them to act on it.

\subsection{From a sudden and unforeseen death}

The Catechism draws the same lesson from the Church's prayers: ``The Church encourages us to prepare ourselves for the hour of our death. In the litany of the saints, for instance, she has us pray: `From a sudden and unforeseen death, deliver us, O Lord'; to ask the Mother of God to intercede for us `at the hour of our death' in the Hail Mary; and to entrust ourselves to St.~Joseph, the patron of a happy death.''\endnote{CCC 1014.} In the Latin of the Litany it runs \latin{A subitanea et improvisa morte}, ``From sudden and unlooked-for death,'' as the Baltimore \emph{Manual of Prayers} renders it.\endnote{\emph{Rituale Romanum} (1872), Litaniae Sanctorum; English from \emph{A Manual of Prayers} (1889), The Litany of the Saints, 260.} Many count a quick and unknowing death a mercy. The Church asks instead for time: time to confess, to receive the last sacraments, to forgive and be forgiven, to make the acts of faith and hope and love, and to say \latin{In manus tuas}. She does not despise those who die suddenly, for Scripture promises, ``But the just man, if he be prevented with death, shall be in rest.''\endnote{Wisdom 4:7, Douay--Rheims (Challoner).}

Each time they say the Hail Mary, the faithful pray, \latin{Sancta Maria, Mater Dei, ora pro nobis peccatoribus, nunc et in hora mortis nostr\ae}: ``Holy Mary, Mother of God, pray for us sinners, now and at the hour of our death.''\endnote{The Hail Mary, Latin and English, \emph{A Manual of Prayers} (1889), 34.} For the Catechism's third counsel, to entrust oneself to Saint Joseph, the Baltimore \emph{Manual} gives the sick three short invocations to say while they can, naming Joseph with Jesus and Mary:

\begin{quote}
Jesus, Mary, and Joseph, I give you my heart and my soul.

Jesus, Mary, and Joseph, assist me in my last agony.

Jesus, Mary, and Joseph, may I breathe forth my soul in peace with you.\endnote{\emph{A Manual of Prayers} (1889), Devotions for the Sick, 470.}
\end{quote}

\subsection{Precious in the sight of the Lord}

Beside the fear of a bad death the saints set the hope of a good one. Alphonsus gives a whole consideration to the death of the just, under the verse ``Precious in the sight of the Lord is the death of his saints.'' ``Considered according to the senses, death excites fear and terror: but viewed with the eyes of faith, it is consoling and desirable. To sinners it appears full of terror; but to the saints it is amiable and precious.'' He borrows Bernard's words: it is precious ``as the end of labours, the consummation of victory, the gate of life.''\endnote{Alphonsus Liguori, \emph{Preparation for Death}, Consideration VIII, first point, Duffy 1843, quoting Bernard, \emph{Sermo in transitu S. Malachiae}.} In his account the just man is not torn from his goods, because his heart was never fastened to them. He does not fear to leave his family, because he loved them in God, and ``at death he recommends them to his heavenly Father, who loves them more than he does.'' He accepts the pains of death with peace, ``and feels consoled in uniting the sacrifice of his death with the sacrifice which Jesus Christ offered for him on the cross to his eternal Father.''\endnote{\emph{Preparation for Death}, Consideration VIII, first point, Duffy 1843.}

Th\'er\`ese of Lisieux died in the Carmel of Lisieux on 30 September 1897, after long suffering. Some weeks before the end she wrote to one of the missionaries she called her brothers: ``I am about to go before God, and I understand now more than ever that one thing only is needful---to work for Him alone, and do nothing for self or creatures. \ldots{} but oh! what joy when comes the happy hour of going Home! I shall not die---I do but enter into Life.''\endnote{\emph{Soeur Th\'er\`ese of Lisieux}, ed. T. N. Taylor (London: Burns, Oates \& Washbourne, 1912; Project Gutenberg text), Letters to Her Brother Missionaries, the letter of 14 August 1897 as Taylor dates it.} On the last evening, when she was told that this was her agony, she answered, ``Ah, very well then\ldots{} very well\ldots{} I do not wish to suffer less!'' Then, looking at her crucifix, she said, ``Oh! \ldots{} I love Him! \ldots{} My God, I\ldots{} love\ldots{} Thee!'' Those were her last words.\endnote{\emph{Soeur Th\'er\`ese of Lisieux}, ed. Taylor, epilogue on her last illness and death.}

Her death is the one the \emph{Craft} describes, given up to God gladly and with a willing heart, and her last words were an act of love.
